% Generated draft: requires mathematical and source review.
\documentclass[11pt]{article}
\usepackage[margin=25mm]{geometry}
\usepackage{fontspec}
\setmainfont{Latin Modern Roman}
\usepackage{amsmath,amssymb,mathtools,unicode-math}
\setmathfont{Latin Modern Math}
\usepackage{enumitem,needspace,xurl,hyperref}
\setlength{\parindent}{0pt}
\setlength{\parskip}{5pt}
\setlength{\emergencystretch}{3em}
\newcommand{\N}{\mathbb{N}}
\newcommand{\Z}{\mathbb{Z}}
\newcommand{\Q}{\mathbb{Q}}
\newcommand{\R}{\mathbb{R}}
\newcommand{\C}{\mathbb{C}}
\newcommand{\F}{\mathbb{F}}
\newcommand{\A}{\mathbb{A}}
\newcommand{\PP}{\mathbb{P}}
\newcommand{\E}{\mathbb{E}}
\newcommand{\eps}{\varepsilon}
\newcommand{\dd}{\,\mathrm{d}}
\begin{document}
\section*{UnsolvedMath ID: 30005048}
\textbf{Source problem code:} catalogue entry
\subsection{Definitions and mathematical statement}
For an interval \(I\subseteq\Z\) and a real sequence \((a_k)_{k\in I-I}\), the associated Toeplitz matrix is \(T=(a_{j-i})_{i,j\in I}\). A matrix is totally nonnegative if every finite square minor has nonnegative determinant. Some literature calls this property ``total positivity''; strict total positivity instead requires every such determinant to be positive. The supplied excerpt does not specify its convention, its index set, or any boundary convention for the sequence.

A common convention defines \(\mathrm{LP}^{+}\) to consist of the nonzero real entire functions that are locally uniform limits on \(\C\) of nonzero real polynomials having nonnegative coefficients and only real, nonpositive zeros. This definition is supplied for orientation: the excerpt does not define \(\mathrm{LP}^{+}\), and its precise convention, including treatment of the zero function and normalization, requires confirmation.

The excerpt attributes to the Aissen--Edrei--Whitney--Schoenberg theorem, with report references [1, 3], the schematic characterization
\[
 T\text{ is totally positive}
 \quad\Longleftrightarrow\quad
 A(t)=\frac{f(t)}{g(-t)},\qquad f,g\in\mathrm{LP}^{+},
\]
where \(A\) is the generating function corresponding to \(T\), and the denominator must not be identically zero. This display records the supplied characterization; it does not assert the equivalence for every index set and sequence in the preceding definition. The excerpt does not specify whether \(A\) is an ordinary or bilateral generating function, its convergence domain, whether the identity is formal or analytic, or the necessary normalization and denominator conditions.

\textbf{Open proof-construction question, as supplied.} Find a more straightforward proof of the intended total-positivity characterization above that replaces the complex-analytic techniques used in report references [3, 4]. The source also explicitly asks about Okounkov's proof [8]. The registry paraphrases the desired replacement as an elementary argument, but neither ``elementary'' nor the permitted proof methods are defined. The exact theorem hypotheses and the intended relation to Okounkov's proof must be recovered from the cited report before this can become a fully specified proof-construction problem. The characterization is presented as an existing theorem; the supplied open question concerns its proof.
\subsection{Short English statement}
The supplied question asks for a more straightforward, described in the registry as elementary, proof of a theorem characterizing totally positive Toeplitz matrices by generating functions of the form \(f(t)/g(-t)\), with \(f,g\in\mathrm{LP}^{+}\). It also mentions Okounkov's proof. The excerpt omits conventions and hypotheses needed to state the theorem completely.
\subsection{Sources}
\begin{itemize}
\item Local UnsolvedMath registry, permanent ID \(30005048\), exact source problem code \texttt{catalogue entry}, title ``Elementary Proofs of Laguerre--P\'olya Total Positivity.'' The supplied record labels the problem open and its statement \texttt{corrected\_verified}; these are supplied curation labels, not independent verification.
\item Supplied source citation: \emph{The Laguerre--P\'olya Class and Combinatorics} (2023), \emph{Oberwolfach Reports}, pp.\ 657--681, DOI \texttt{10.4171/owr/2022/13}. Supplied source and literature-source link: \url{https://doi.org/10.4171/owr/2022/13}. The record describes Oberwolfach Reports as documenting mathematical workshops and their open-problem discussions.
\item The excerpt attributes ``Total-positivity characterization of the Laguerre--P\'olya class \(\mathrm{LP}^{+}\)'' to Alexander Dyachenko, joint work with Alan Sokal. It names the Aissen--Edrei--Whitney--Schoenberg theorem with report references [1, 3], complex-analytic techniques in [3, 4], and Okounkov's proof [8]. Bibliographic entries for those numbered references were not supplied.
\item Literature assessment checked at \texttt{2026-08-22}: the supplied curation says that the primary source explicitly presents replacement of the complex-analytic proofs, including discussion of Okounkov's proof, as an intriguing open question, and that no later resolution was located. The listed source checked was the Oberwolfach report linked above. This is supplied dated curation, not an independently verified literature assessment; the record requests independent verification of time-sensitive status and cited claims.
\item The record gives proposed year 2022 and states that its clean formulation was checked against the official report. These provenance statements are supplied metadata; no source was visited for this draft.
\end{itemize}
\end{document}
